\documentclass[11pt]{ctexart}
\usepackage[a4paper,margin=1in]{geometry}
\usepackage{graphicx}
\usepackage{xcolor}
\usepackage{hyperref}
\definecolor{refgray}{gray}{0.45}
\title{\textbf{市场最新观点汇总}\\[0.4em]{\large 能源供给冲击与AI需求共振，通胀韧性重塑全球资产定价逻辑}}
\author{KC桌面 自动生成}
\date{260517}
\begin{document}
\maketitle
\section*{一页摘要}
\begin{itemize}
  \item 全球市场正从“定价尾部风险消失”转向“定价短缺延续”，能源供应链中断与AI基础设施瓶颈成为资本追逐的主线。
  \item 通胀风险溢价在全球收益率曲线上系统性重定价，美债30年期收益率突破5\%并非孤立事件，而是全球长端利率同步上行的缩影。
  \item 存储芯片市场经历结构性变化，NAND定价权从规模最大的供应商向技术领先者转移，供给紧张预计持续到2027年。
  \item 中国经济呈现“地产暖、出行冷”的分化格局，能源成本冲击与房地产边际企稳之间的张力成为核心矛盾。
  \item 欧洲经济面临供给侧滞后压力，能源价格向核心通胀的传导呈现“火箭与羽毛”式不对称，市场对通胀风险分布存在系统性偏差。
  \item 央行购金进入“不可见”新阶段，伦敦金库流出与英国海关出口数据之间的缺口累计约250吨，结构性需求强度被低估。
  \item 全球资金流向显示，市场正在从“通胀交易”的短期博弈切换到“通胀韧性下的结构性受益者”的长期布局，科技与基础设施基金持续吸金。
\end{itemize}
\section{宏观与利率：通胀风险溢价重定价，全球长端利率同步上行}
\textbf{市场正在经历的不只是通胀预期的上修，而是通胀风险溢价在整个收益率曲线上的系统性重定价，这将从根本上改变债券作为组合对冲工具的有效性。}

\begin{itemize}
  \item 美债30年期收益率突破5\%并非孤立事件，而是全球长端利率同步上行的一部分，英国、日本和德国30年期收益率均值已上升约50个基点。
  \item 通胀风险溢价上升意味着市场开始为“通胀路径的不确定性”定价，这种不确定性无法被单一的加息或降息所消除。
  \item 能源价格向核心通胀的传导呈现“火箭与羽毛”式不对称：能源涨价时传导效率超过100\%，而降价时仅约58\%，导致核心通胀更加顽固。
  \item 欧洲经济正在被重新定价，但市场低估了供给侧冲击的滞后效应，能源价格对GDP的拖累约0.25个百分点，欧元区2026年GDP增速预期已从1.4\%下调至0.5\%。
  \item 菲律宾正经历滞胀式能源供给冲击，4月CPI升至7.2\%，一季度GDP增速骤降至2.8\%，政治尾部风险对投资的压制被市场低估。
\end{itemize}
\begin{figure}[htbp]
\centering
\includegraphics[width=0.88\linewidth]{figures/fig_066.jpg}
\caption{Exhibit 2 - Exhibit 1: 30y UST yields' renewed break of \$5\textbackslash{}\%\$ has come amid a sustained uptrend in global long-end yields UST 30y yield vs avg of UK, German, Japan 30y yields}
\end{figure}
\begin{figure}[htbp]
\centering
\includegraphics[width=0.88\linewidth]{figures/fig_067.jpg}
\caption{Exhibit 2 - Exhibit 2: Treasuries have been a poor portfolio diversifier since the start of the conflict}
\end{figure}
\begin{figure}[htbp]
\centering
\includegraphics[width=0.88\linewidth]{figures/fig_069.jpg}
\caption{Exhibit 4 - Exhibit 4: Inflation forwards are no longer clearly lagging the change in oil prices}
\end{figure}
\begin{figure}[htbp]
\centering
\includegraphics[width=0.88\linewidth]{figures/fig_070.jpg}
\caption{Exhibit 5 - Exhibit 5: Forward HICP pricing reasonable degree of passthrough from headline to core inflation Change in 2y2y HICP and model fitted value since 1 Jan.}
\end{figure}
\begin{figure}[htbp]
\centering
\includegraphics[width=0.88\linewidth]{figures/fig_061.jpg}
\caption{Exhibit 1 - Exhibit 1: Previous Studies Point to Uneven Passthrough of Energy Shocks}
\end{figure}
\begin{figure}[htbp]
\centering
\includegraphics[width=0.88\linewidth]{figures/fig_062.jpg}
\caption{Exhibit 2 - Exhibit 2: Top-Down Evidence Clearly Shows an Asymmetry in Food and Energy Shock Passthrough}
\end{figure}
\begin{figure}[htbp]
\centering
\includegraphics[width=0.88\linewidth]{figures/fig_065.jpg}
\caption{Exhibit 5 - Exhibit 5: The Pricing Asymmetry Creates Upside Risks for Core Inflation}
\end{figure}
\begin{figure}[htbp]
\centering
\includegraphics[width=0.88\linewidth]{figures/fig_053.jpg}
\caption{Exhibit 1 - Exhibit 1: Current Shock is Seen as More Contained for European Natural Gas Than in 2022; Oil Shock Expected to Be More Persistent}
\end{figure}
\begin{figure}[htbp]
\centering
\includegraphics[width=0.88\linewidth]{figures/fig_057.jpg}
\caption{Exhibit 5 - Exhibit 5: Our Model Points to a 0.25pp Drag on Growth This Year from the Energy Shock}
\end{figure}
\begin{figure}[htbp]
\centering
\includegraphics[width=0.88\linewidth]{figures/fig_058.jpg}
\caption{Exhibit 6 - Exhibit 6: We Have Downgraded GDP Growth in 2026 by Around 0.8pp in the Euro Area; Pre-War Data Surprised to the Upside in the UK}
\end{figure}
{\small\color{refgray}\textbf{References:} [R005] 欧洲经济正在被重新定价，但市场低估了供给侧的滞后压力; [R006] 市场真正低估的不是通胀读数，而是能源价格传导的“不对称性”; [R007] 市场真正低估的不是通胀，而是通胀风险溢价的重定价; [R002] 菲律宾正经历一场滞胀式能源供给冲击，市场低估了政治尾部风险对增长的压制}

\section{AI/算力与存储：NAND供给侧再定价，光通信市场增速大幅上调}
\textbf{市场真正低估的不是AI需求，而是NAND供给侧的再定价，定价权正在从买家向卖家转移；光通信市场增速被大幅上调，但CPO的爆发要等到2028年之后。}

\begin{itemize}
  \item Kioxia给出CY2Q26收入环比增长70\%的指引，几乎全部来自ASP提升，暗示市场对三星和SK海力士的NAND ASP预测可能过于保守。
  \item NAND供给紧张预计持续到2027年，三家主要供应商的定价行为趋于同步，行业盈利天花板尚未到来。
  \item 三星潜在的劳资纠纷可能成为定价催化剂，后端封测环节的劳动密集性使其更易受影响。
  \item LightCounting将Datacom市场2025-2028年CAGR上调至近+35\%，1.6T是增长引擎，800G则是稳定的“压舱石”。
  \item 每颗XPU的光模块用量持续爬升，从2023年的2.7个预计增至2027年的接近4.5个，但不同架构差异显著。
  \item CPO市场被低估，第三方预测2028年达20亿美元，但Lumentum自身指引显示仅一家公司2028年CPO收入就可能达2亿美元。
\end{itemize}
\begin{figure}[htbp]
\centering
\includegraphics[width=0.88\linewidth]{figures/fig_007.jpg}
\caption{Exhibit 11 - Exhibit 11: NAND wafer contract price rally significantly higher vs spot NAND wafer spot and contract price - quarterly average trend}
\end{figure}
\begin{figure}[htbp]
\centering
\includegraphics[width=0.88\linewidth]{figures/fig_015.jpg}
\caption{Exhibit 20 - Exhibit 20: Current contract price at \textbackslash{}\textasciitilde{}\textbackslash{}\$25, around 10x vs Feb-25 bottom of \textbackslash{}\$2.5; prices significantly up in 4Q25 and through 1Q26 NAND wafer contract price trend, Jun '23-Apr '26}
\end{figure}
\begin{figure}[htbp]
\centering
\includegraphics[width=0.88\linewidth]{figures/fig_020.jpg}
\caption{Exhibit 24 - Exhibit 24: Current NAND spot and contract prices are a few times higher than 2025 summer level 512Gb NAND wafer spot vs contract price trend (US\textbackslash{}\$), Apr '19 - Apr '26}
\end{figure}
\begin{figure}[htbp]
\centering
\includegraphics[width=0.88\linewidth]{figures/fig_026.jpg}
\caption{Exhibit 31 - Exhibit 31: Robust rally continues led by NAND/HDD companies, followed by DRAM names, memory companies share price significantly up this week driven by AI-led demand Memory companies – 2026 YTD stock performance comparison}
\end{figure}
\begin{figure}[htbp]
\centering
\includegraphics[width=0.88\linewidth]{figures/fig_285.jpg}
\caption{Figure 1 - Figure 1: Datacom Market Forecast by Data Rate \textbackslash{}\$ in Billions}
\end{figure}
\begin{figure}[htbp]
\centering
\includegraphics[width=0.88\linewidth]{figures/fig_286.jpg}
\caption{Figure 2 - Figure 2: 800G Datacom Market Forecast \textbackslash{}\$ in Billions}
\end{figure}
\begin{figure}[htbp]
\centering
\includegraphics[width=0.88\linewidth]{figures/fig_287.jpg}
\caption{Figure 3 - Figure 3: 1.6T Datacom Market Forecast \textbackslash{}\$ in Billions}
\end{figure}
\begin{figure}[htbp]
\centering
\includegraphics[width=0.88\linewidth]{figures/fig_288.jpg}
\caption{Figure 4 - Figure 4: Datacom Attach Rate X}
\end{figure}
\begin{figure}[htbp]
\centering
\includegraphics[width=0.88\linewidth]{figures/fig_290.jpg}
\caption{Figure 6 - Figure 6: Long-term CPO Market Forecast \textbackslash{}\$ in Billions}
\end{figure}
{\small\color{refgray}\textbf{References:} [R001] 存储芯片的下一轮定价权，不在三星手里; [R015] 市场真正低估的不是AI需求，而是NAND供给侧的再定价; [R020] 光通信市场真正的转折点不在2025，而在2028年之后}

\section{能源与大宗：供给纪律兑现驱动钢价上涨，黄金购金进入“不可见”阶段}
\textbf{全球钢铁市场的真正拐点不在需求侧，而在供给纪律的兑现；央行购金正在脱离传统可观测渠道，结构性需求强度被低估。}

\begin{itemize}
  \item 全球钢价上涨发生在粗钢产量同比收缩的背景下，巴西热轧卷板价格年初至今累计涨幅超过20\%，美国涨14.7\%。
  \item 区域价格分化的背后是产能利用率的结构性差异，除中国外地区有效产能利用率达99\%，印度持续满产。
  \item 中国5月前两周粗钢产量同比下降3.2\%，但“去内卷”和长期产能削减计划的执行力度可能延迟。
  \item 自2025年8月起，伦敦金库流出与英国海关出口数据之间出现系统性缺口，累计“失踪”主权黄金流量约240吨。
  \item 修正后的央行购金模型显示月均购金量从约29吨上调至约50吨/月，主权购金行为透明度下降。
\end{itemize}
\begin{figure}[htbp]
\centering
\includegraphics[width=0.88\linewidth]{figures/fig_097.jpg}
\caption{Exhibit 1 - Exhibit 1: HRC Prices in major regions \textbackslash{}\$/t}
\end{figure}
\begin{figure}[htbp]
\centering
\includegraphics[width=0.88\linewidth]{figures/fig_098.jpg}
\caption{Exhibit 2 - Exhibit 2: HRC 1m/YTD price performance, \textbackslash{}\$ terms Data for April 2026 Exhibit 3: HRC price ranges in the past 4 years US\textbackslash{}\$/t}
\end{figure}
\begin{figure}[htbp]
\centering
\includegraphics[width=0.88\linewidth]{figures/fig_099.jpg}
\caption{Exhibit 4 - Exhibit 4: Rebar prices in major regions \textbackslash{}\$/t}
\end{figure}
\begin{figure}[htbp]
\centering
\includegraphics[width=0.88\linewidth]{figures/fig_103.jpg}
\caption{Exhibit 9 - Exhibit 9: Effective steel capacity utilization Latest monthly output as \% of peak output (since 2022)}
\end{figure}
\begin{figure}[htbp]
\centering
\includegraphics[width=0.88\linewidth]{figures/fig_104.jpg}
\caption{Exhibit 10 - Exhibit 10: Ex-China effective steel capacity utilization Latest monthly output as \% of peak output (since 2022)}
\end{figure}
\begin{figure}[htbp]
\centering
\includegraphics[width=0.88\linewidth]{figures/fig_107.jpg}
\caption{Exhibit 13 - Exhibit 13: China crude steel output change \% m/m vs 5Y median trend}
\end{figure}
\begin{figure}[htbp]
\centering
\includegraphics[width=0.88\linewidth]{figures/fig_174.jpg}
\caption{Exhibit 2 - Exhibit 1: We Update Our Nowcast of Central Bank Gold Purchases, Because...}
\end{figure}
\begin{figure}[htbp]
\centering
\includegraphics[width=0.88\linewidth]{figures/fig_175.jpg}
\caption{Exhibit 2 - Exhibit 2: ... A Consistent Gap Between London Vault Outlfows and UK Net Exports Suggests Our GS Nowcast Had Been Undershooting Central Bank Gold Purchases Since August 2025}
\end{figure}
{\small\color{refgray}\textbf{References:} [R010] 全球钢铁市场的真正拐点，不在需求侧，而在供给纪律的兑现; [R012] 市场尚未充分定价的，是央行购金正在进入一个“不可见”的新阶段}

\section{地产与消费：中国房地产边际企稳，消费出行降温}
\textbf{中国经济最值得关注的是两个结构性背离——能源成本冲击与房地产边际企稳之间的张力，以及消费信心修复与实物消费量不匹配之间的裂缝。}

\begin{itemize}
  \item 30城新房日均成交量回升且高于去年同期，16城二手房成交量也出现反弹，房地产市场可能进入“基数效应驱动下的边际企稳”阶段。
  \item 二手房成交量反弹幅度更大，说明市场仍在“以价换量”，新房市场回暖更多依赖国企开发商和核心城市供给推动。
  \item 国内客运航班量大幅下降，航班取消率飙升，交通拥堵指数上升，表明短途出行增多、长途减少。
  \item 消费者信心指数升至2019年以来最高，但实际消费行为偏谨慎，4月汽车销量低于去年同期。
  \item 能源价格冲击正在传导，5月11日国内汽柴油价格分别上调320元/吨和310元/吨，4月原油进口价格环比大幅上涨。
\end{itemize}
\begin{figure}[htbp]
\centering
\includegraphics[width=0.88\linewidth]{figures/fig_030.jpg}
\caption{Exhibit 1 - Exhibit 1: 30-city daily property transaction volume in the primary market increased over the last week and was above year-ago level}
\end{figure}
\begin{figure}[htbp]
\centering
\includegraphics[width=0.88\linewidth]{figures/fig_031.jpg}
\caption{Exhibit 2 - Exhibit 2: 16-city daily property transaction volume in the secondary market rebounded and was above year-ago level}
\end{figure}
\begin{figure}[htbp]
\centering
\includegraphics[width=0.88\linewidth]{figures/fig_032.jpg}
\caption{Exhibit 3 - Exhibit 3: China's passenger flights for domestic routes declined sharply over the past week, and were below year ago level}
\end{figure}
\begin{figure}[htbp]
\centering
\includegraphics[width=0.88\linewidth]{figures/fig_033.jpg}
\caption{Exhibit 4 - Exhibit 4: China's passenger flights cancellation rate surged over the last week, and was above year-ago level}
\end{figure}
\begin{figure}[htbp]
\centering
\includegraphics[width=0.88\linewidth]{figures/fig_034.jpg}
\caption{Exhibit 5 - Exhibit 5: Traffic congestion increased over the last week}
\end{figure}
\begin{figure}[htbp]
\centering
\includegraphics[width=0.88\linewidth]{figures/fig_035.jpg}
\caption{Exhibit 6 - Exhibit 6: Domestic gasoline and diesel prices were raised by 320 and 310 RMB/tonne respectively on 11 May}
\end{figure}
\begin{figure}[htbp]
\centering
\includegraphics[width=0.88\linewidth]{figures/fig_036.jpg}
\caption{Exhibit 8 - Exhibit 8: Import prices for crude oil and refined petroleum products surged in April}
\end{figure}
\begin{figure}[htbp]
\centering
\includegraphics[width=0.88\linewidth]{figures/fig_038.jpg}
\caption{Exhibit 10 - Exhibit 10: The Morning Consult consumer confidence rose to its highest level since 2019}
\end{figure}
\begin{figure}[htbp]
\centering
\includegraphics[width=0.88\linewidth]{figures/fig_039.jpg}
\caption{Exhibit 11 - Exhibit 11: New energy vehicles (NEVs) sales volume edged up in April, but remained below year-ago level}
\end{figure}
\begin{figure}[htbp]
\centering
\includegraphics[width=0.88\linewidth]{figures/fig_040.jpg}
\caption{Exhibit 12 - Exhibit 12: Total auto sales volume declined in April and remained below year-ago level}
\end{figure}
{\small\color{refgray}\textbf{References:} [R004] 经济活动的真相藏在能源和房地产的背离里}

\section{区域市场：日本动量拥挤，韩国宽度狭窄，印度资金再分配}
\textbf{日本市场正经历定价逻辑的结构性切换，动量交易已拥挤到对好消息完全免疫；韩国指数大涨背后市场宽度创纪录狭窄；印度市场正经历由宏观冲击驱动的资金再分配。}

\begin{itemize}
  \item 日本财报季54\%的公司实现正向意外，但动量股在出现负向意外后3个交易日内中位数相对收益跑输TOPIX约5个百分点，而对正向意外几乎无反应。
  \item 日本“关键资源”篮子年初至今涨幅43\%，是TOPIX同期涨幅的三倍有余，受益于“经济安全”叙事。
  \item KOSPI年初至今涨幅近8成，但835只上市股票中仅68只跑赢指数，这些股票占据近66\%的市值。
  \item 韩国市场宽度指标从-1标准差下方反弹，大盘股相对中小盘的超额收益已超过2个标准差，均值回归力量可能正在积蓄。
  \item 印度本土机构投资者与外资预期脱钩，外资在撤、本土资金在接，但接盘的逻辑和定价权正在发生转移。
  \item 印度4月出口同比增长13.8\%，电子产品出口增速达40.3\%，出口结构正从传统劳动密集型向电子制造转变。
\end{itemize}
\begin{figure}[htbp]
\centering
\includegraphics[width=0.88\linewidth]{figures/fig_225.jpg}
\caption{Exhibit 1 - Exhibit 1: GSJPCRTT and GSJPCRTR have outperformed TOPIX Data as of May 14, 2026}
\end{figure}
\begin{figure}[htbp]
\centering
\includegraphics[width=0.88\linewidth]{figures/fig_226.jpg}
\caption{Exhibit 2 - Exhibit 2: So far, 4Q positive surprises exceed negative surprises TOPIX constituent companies with Feb/Mar fiscal year-end, as of May 14}
\end{figure}
\begin{figure}[htbp]
\centering
\includegraphics[width=0.88\linewidth]{figures/fig_230.jpg}
\caption{Exhibit 8 - Exhibit 8: Momentum stocks reacted strongly to negative surprises and did not react to positive surprises 1M momentum stocks with Mar fiscal year end, median TOPIX-relative price reaction, as of May 15}
\end{figure}
\begin{figure}[htbp]
\centering
\includegraphics[width=0.88\linewidth]{figures/fig_231.jpg}
\caption{Exhibit 9 - Exhibit 9: Spike in standard deviation of +1D price reactions As of May 15}
\end{figure}
\begin{figure}[htbp]
\centering
\includegraphics[width=0.88\linewidth]{figures/fig_253.jpg}
\caption{Exhibit 8 - Exhibit 1: While Korea's daily market breadth has remained narrow, it has recently rebounded from below the -1 standard deviation level}
\end{figure}
\begin{figure}[htbp]
\centering
\includegraphics[width=0.88\linewidth]{figures/fig_254.jpg}
\caption{Exhibit 3 - Exhibit 3: Only 68 out of 835 listed stocks have outperformed the index, and these names account for nearly 66\% of total market capitalization}
\end{figure}
\begin{figure}[htbp]
\centering
\includegraphics[width=0.88\linewidth]{figures/fig_255.jpg}
\caption{Exhibit 2 - Exhibit 2: The KOSPI 50 (large-cap index) has outperformed the KOSPI 200 small- and mid-cap segment by more than 2 standard deviations}
\end{figure}
\begin{figure}[htbp]
\centering
\includegraphics[width=0.88\linewidth]{figures/fig_256.jpg}
\caption{Exhibit 4 - Exhibit 4: Despite the overall market rally, nearly 70\% of KOSPI-listed stocks are still trading below book value, which we believe presents attractive valuation opportunities}
\end{figure}
\begin{figure}[htbp]
\centering
\includegraphics[width=0.88\linewidth]{figures/fig_167.jpg}
\caption{Exhibit 1 - Exhibit 1: Demand continues to outpace supply for DRAM and NAND, leading to exceptionally strong memory producer profitability Demand / Supply Balance for DRAM and NAND (\%)}
\end{figure}
\begin{figure}[htbp]
\centering
\includegraphics[width=0.88\linewidth]{figures/fig_170.jpg}
\caption{Exhibit 4 - Exhibit 4: Historically, we have seen India earnings get downgraded by 6-13\% (median 7\%) in the 12 months after an energy supply shock}
\end{figure}
\begin{figure}[htbp]
\centering
\includegraphics[width=0.88\linewidth]{figures/fig_171.jpg}
\caption{Exhibit 5 - Exhibit 5: After continued outflows from active global money, domestic clients expressed worry around tapering passive inflows}
\end{figure}
{\small\color{refgray}\textbf{References:} [R018] 日本市场的真正分化，不在行业之间，而在“动量”与“价值”的定价权转移; [R019] 韩国市场真正的机会，藏在指数大涨背后的结构性扭曲里; [R011] 印度市场的真正分歧不在涨跌，而在“谁先扛不住”和“谁在悄悄换仓”; [R014] 印度贸易逆差扩大的真正信号：出口韧性背后，是结构性再定价的开始}

\section{风险偏好与资金流向：市场为“通胀韧性”重新定价资产}
\textbf{当前资金流动的主导逻辑并非对经济衰退担忧的消退，而是对“通胀中枢上移”这一结构性事实的逐步定价，市场正在筛选“通胀韧性”的真正受益者。}

\begin{itemize}
  \item 截至5月13日当周，全球股票基金净流入200亿美元，固定收益基金净流入285亿美元，风险偏好全面修复。
  \item 技术基金和基础设施基金录得全行业最大净流入，消费品基金遭遇最大净流出，能源基金在连续大额流入后转为净流出5亿美元。
  \item 欧洲财报季EPS超预期约2.3\%，但市场对好消息的奖励正在失效，超预期的公司仅获得不到1\%的超额收益。
  \item 欧洲资本品板块订单势头强劲（77\%的公司超预期），但订单转化为销售的时间在拉长，约50\%的公司利润率承压。
  \item 全球汇率市场的核心驱动力正在从利差和风险偏好转向贸易条件，能源进口依赖度高且缺乏高技术出口对冲的经济体货币承压。
\end{itemize}
\begin{figure}[htbp]
\centering
\includegraphics[width=0.88\linewidth]{figures/fig_176.jpg}
\caption{Exhibit 1 - Exhibit 1: Net Unhedged Flows into US Equity Funds}
\end{figure}
\begin{figure}[htbp]
\centering
\includegraphics[width=0.88\linewidth]{figures/fig_177.jpg}
\caption{Exhibit 3 - Exhibit 1: EPS has surprised by around 2\%, only slightly below the historical average Equal Weighted (\%)}
\end{figure}
\begin{figure}[htbp]
\centering
\includegraphics[width=0.88\linewidth]{figures/fig_178.jpg}
\caption{Exhibit 3 - Exhibit 3: Guidance surprises are seeing small price reactions SXXP companies. Average price reaction on reporting day (vs. the market). FY1 guidance (vs. consensus one day before the guidance was issued).}
\end{figure}
\begin{figure}[htbp]
\centering
\includegraphics[width=0.88\linewidth]{figures/fig_179.jpg}
\caption{Exhibit 2 - Exhibit 2: Market reactions remain asymmetric Average price reaction on reporting day vs. the market - STOXX 600 companies}
\end{figure}
\begin{figure}[htbp]
\centering
\includegraphics[width=0.88\linewidth]{figures/fig_180.jpg}
\caption{Exhibit 4 - Exhibit 4: Energy and Basic Resources names continue to guide the aggregate upgrades 2026 EPS revisions. STOXX 600 sectors.}
\end{figure}
\begin{figure}[htbp]
\centering
\includegraphics[width=0.88\linewidth]{figures/fig_073.jpg}
\caption{Exhibit 1 - GS \& Co. LLC}
\end{figure}
\begin{figure}[htbp]
\centering
\includegraphics[width=0.88\linewidth]{figures/fig_074.jpg}
\caption{Exhibit 2 - Exhibit 2: ...while inflation projections have continued to drift higher}
\end{figure}
\begin{figure}[htbp]
\centering
\includegraphics[width=0.88\linewidth]{figures/fig_075.jpg}
\caption{Exhibit 3 - Exhibit 3: Recent divergence between actual and model-implied EUR/GBP demonstrates a partial reintroduction of political premium in Sterling}
\end{figure}
{\small\color{refgray}\textbf{References:} [R013] 市场正在为“通胀韧性”重新定价资产，而非衰退风险; [R016] 欧洲财报季的真正信号：市场对好消息的奖励正在失效; [R008] 市场真正低估的不是美元强弱，而是贸易条件正在重塑全球汇率的定价锚}

\section*{结语}
综合来看，当前全球市场的核心主线是“短缺”的再定价——无论是能源、存储芯片、AI基础设施还是黄金，供给侧的制约正在成为资产定价的主导力量。通胀韧性取代衰退担忧成为资金配置的主逻辑，但市场对通胀风险分布、供给冲击粘性以及结构性资金流动的定价仍不充分。理解这些被低估的变量，比判断短期涨跌更为关键。
\vfill
{\small\color{refgray}Personal reading notes and learning share only. Not investment advice.}
\end{document}
